% OpenStax Calculus Volume 2 -- Section 4.5 Exercises: First-Order Linear Equations
% Exercises reproduced from OpenStax, licensed under CC BY 4.0.
% Source: https://openstax.org/books/calculus-volume-2/pages/4-5-first-order-linear-equations
\documentclass{exercises}
\title{OpenStax Calculus Volume 2}
\subtitle{Section 4.5 Exercises: First-Order Linear Equations}
\sourceurl{https://openstax.org/books/calculus-volume-2/pages/4-5-first-order-linear-equations}
\author{}
\date{}

\begin{document}
\maketitle


Are the following differential equations linear? Explain your reasoning.

\begin{enumerate}[start=1]
\item $\frac{dy}{dx}=x^{2}y+\sin x$
\item $\frac{dy}{dt}=ty$
\item $\frac{dy}{dt}+y^{2}=t$
\item $y'=x^{3}+e^{x}$
\item $y'=y+e^{y}$
\end{enumerate}

Write the following first-order differential equations in standard form.

\begin{enumerate}[resume]
\item $y'=x^{3}y+\sin x$
\item $y'+3y-\ln x=0$
\item $-xy'=(3x+2)y+xe^{x}$
\item $\frac{dy}{dt}=4y+ty+\tan t$
\item $\frac{dy}{dt}=yt(t+1)$
\end{enumerate}

What are the integrating factors for the following differential equations?

\begin{enumerate}[resume]
\item $y'=xy+3$
\item $y'+e^{x}y=\sin x$
\item $y'=x\ln(x)y+3x$
\item $\frac{dy}{dx}=\tanh(x)y+1$
\item $\frac{dy}{dt}+3ty=e^{t}y$
\end{enumerate}

Solve the following differential equations by using integrating factors.

\begin{enumerate}[resume]
\item $y'=3y+2$
\item $y'=2y-x^{2}$
\item $xy'=3y-6x^{2}$
\item $(x+2)y'=3x+y$
\item $y'=3x+xy$
\item $xy'=x+y$
\item $\sin(x)y'=y+2x$
\item $y'=y+e^{x}$
\item $xy'=3y+x^{2}$
\item $y'+\ln x=\frac{y}{x}$
\end{enumerate}

Solve the following differential equations. Use your calculator to draw a family of solutions. Are there certain initial conditions that change the behavior of the solution?

\begin{enumerate}[resume]
\item \techrequired $(x+2)y'=2y-1$
\item \techrequired $y'=3e^{t/3}-2y$
\item \techrequired $ty'+\frac{y}{2}=\sin(3t)$
\item \techrequired $xy'=2\frac{\cos x}{x}-3y$
\item \techrequired $(x+1)y'=3y+x^{2}+2x+1$
\item \techrequired $\sin(x)y'+\cos(x)y=2x$
\item \techrequired $\sqrt{x^{2}+1}y'=y+2$
\item \techrequired $x^{3}y'+2x^{2}y=x+1$
\end{enumerate}

Solve the following initial-value problems by using integrating factors.

\begin{enumerate}[resume]
\item $y'+y=x$, $y(0)=3$
\item $y'=y+2x^{2}$, $y(0)=0$
\item $xy'=y-3x^{3}$, $y(1)=0$
\item $x^{2}y'=xy-\ln x$, $y(1)=1$
\item $(1+x^{2})y'=y-1$, $y(0)=0$
\item $xy'=y+2x\ln x$, $y(1)=5$
\item $(2+x)y'=y+2+x$, $y(0)=0$
\item $y'=xy+2xe^{\frac{1}{2}x^{2}}$, $y(0)=2$
\item $\sqrt{x}y'=y+2x$, $y(0)=1$
\item $y'=2y+xe^{x}$, $y(0)=-1$
\item A falling object of mass $m$ can reach terminal velocity when the drag force is proportional to its velocity, with proportionality constant $k$. Set up the differential equation and solve for the velocity given an initial velocity of $0$.
\item Using your expression from the preceding problem, what is the terminal velocity? (Hint: Examine the limiting behavior; does the velocity approach a value?)
\item \techrequired Using your equation for terminal velocity, solve for the distance fallen. How long does it take to fall $5000$ meters if the mass is $100$ kilograms, the acceleration due to gravity is $9.8\,\text{m}/\text{s}^{2}$ and the proportionality constant is $4$?
\item A more accurate way to describe terminal velocity is that the drag force is proportional to the square of velocity, with a proportionality constant $k$. Set up the differential equation and solve for the velocity.
\item Using your expression from the preceding problem, what is the terminal velocity? (Hint: Examine the limiting behavior: Does the velocity approach a value?)
\item \techrequired Using your equation for terminal velocity, solve for the distance fallen. How long does it take to fall $5000$ meters if the mass is $100$ kilograms, the acceleration due to gravity is $9.8\,\text{m}/\text{s}^{2}$ and the proportionality constant is $4$? Does it take more or less time than your initial estimate?
\end{enumerate}

For the following problems, determine how parameter $a$ affects the solution.

\begin{enumerate}[resume]
\item Solve the generic equation $y'=ax+y$. How does varying $a$ change the behavior?
\item Solve the generic equation $y'=ay+x$. How does varying $a$ change the behavior?
\item Solve the generic equation $y'=ax+xy$. How does varying $a$ change the behavior?
\item Solve the generic equation $y'=x+axy$. How does varying $a$ change the behavior?
\item Solve $y'-y=e^{kt}$ with the initial condition $y(0)=0$. As $k$ approaches $1$, what happens to your formula?
\end{enumerate}

\end{document}
